% Generated draft: requires mathematical and source review.
\documentclass[11pt]{article}
\usepackage[margin=25mm]{geometry}
\usepackage{fontspec}
\setmainfont{Latin Modern Roman}
\usepackage{amsmath,amssymb,mathtools,unicode-math}
\setmathfont{Latin Modern Math}
\usepackage{enumitem,needspace,xurl,hyperref}
\setlength{\parindent}{0pt}
\setlength{\parskip}{5pt}
\setlength{\emergencystretch}{3em}
\newcommand{\N}{\mathbb{N}}
\newcommand{\Z}{\mathbb{Z}}
\newcommand{\Q}{\mathbb{Q}}
\newcommand{\R}{\mathbb{R}}
\newcommand{\C}{\mathbb{C}}
\newcommand{\F}{\mathbb{F}}
\newcommand{\A}{\mathbb{A}}
\newcommand{\PP}{\mathbb{P}}
\newcommand{\E}{\mathbb{E}}
\newcommand{\eps}{\varepsilon}
\newcommand{\dd}{\,\mathrm{d}}
\begin{document}
\section*{UnsolvedMath ID: 801064}
\textbf{Source problem code:} catalogue entry
\subsection{Definitions and mathematical statement}
The supplied record does not define the map \(H_2\), give its formula, or specify its parameters. The following question therefore requires identification of that map, with the displayed iterates defined on the relevant points.

For \(h>0\), let
\[
V_h=\{z\in\C/\Z:\operatorname{Im}z<h\},
\]
where \(\C/\Z\) is the quotient by integer translations. For an irrational rotation class \(\theta\in(\R\setminus\Q)/\Z\), translation \(z\mapsto z+\theta\) is well-defined on \(V_h\). Following the supplied definition, say that \(H_2\) has a Herman ring with rotation number \(\theta\) if there exist \(h>0\), a holomorphic map \(\phi:V_h\to\C^2\), and an integer \(n\geq2\) such that
\[
H_2^{\circ n}(\phi(z))=\phi(z+\theta)
\qquad\text{for every }z\in V_h,
\]
where \(H_2^{\circ n}\) denotes the \(n\)-fold iterate.

\textbf{Question.} Does the map denoted by \(H_2\) admit such data for at least one \(\theta\in(\R\setminus\Q)/\Z\)?

The domain above retains exactly the supplied one-sided inequality; no lower bound on \(\operatorname{Im}z\) is imposed. The supplied definition requires only that \(\phi\) be holomorphic, without requiring it to be nonconstant, injective, or an embedding. These omissions leave the intended geometric meaning of ``Herman ring'' unresolved; no additional conditions have been imposed here.
\subsection{Short English statement}
For the unspecified map \(H_2\), ask whether some iterate of order \(n\geq2\) intertwines a holomorphic map from the stated annulus with an irrational translation. The map \(H_2\) and the intended nondegeneracy requirements on the holomorphic map need clarification.
\subsection{Sources}
\begin{itemize}
\item Local UnsolvedMath registry: permanent ID \(801064\), source problem code \texttt{catalogue entry}, title ``Henon Maps: Open Problems --- Question 16.'' The supplied record attributes the problem to Berger, Bedford, with proposed year 2023, and labels its status as open; this status is supplied metadata, not an independently verified assessment.
\item Berger, Bedford, \emph{Henon Maps: A List of Open Problems} (2023), Question 16. Supplied source URL: \url{https://arxiv.org/abs/2312.03907}. The record reports access on 2026-10-16 and extraction by \texttt{arxiv\_latex}; these are supplied provenance statements, not a claim of independent source access or verification.
\item Supplied upstream revision: \texttt{d371ee917ac9647fbb06cd6fc529deec6f3dda82}.
\end{itemize}
\end{document}
